\documentclass[11pt,ngerman]{scrartcl}
\usepackage[T1]{fontenc}
\usepackage[utf8]{inputenc}
\usepackage[ngerman]{babel}
\usepackage[a4paper,margin=2.5cm]{geometry}
\usepackage[scaled]{helvet}
\renewcommand{\familydefault}{\sfdefault}
\usepackage{enumitem}
\setlist[itemize]{leftmargin=1.5em,noitemsep,nolistsep}

\begin{document}

\begin{center}
    {\LARGE \textbf{Rundgang}}
\end{center}
\vspace{-1em}

\subsection*{Mensa}
\begin{itemize}
    \item Karte aufladen
    \item Prinzip von Theken
    \item Prinzip vom Bezahlen
\end{itemize}

\subsection*{Empfang}
\begin{itemize}
    \item Schlüssel für Räume ausleihen
\end{itemize}

\subsection*{Cafeteria}

\subsection*{Hörsääle 1-8}

\subsection*{Fachschaft}
\begin{itemize}
    \item 3D-Druck
    \item Entspannen
    \item Auf Ersti-Heft aufmerksam machen
    \item Ansprechpartner:Innen / Vertreter:Innen für Studierendenschaft
\end{itemize}

\subsection*{Zentralausleihe/Service-Point A041.1}
\begin{itemize}
    \item Raumausstattung (Kreide, Folien, Beamer, Laptops, etc...)
    \item Elektronik Zeug ausleihen (Microcontroller, Widerstände, Dioden, FPGAs...)
\end{itemize}

\subsection*{Maschinenhalle}
\begin{itemize}
    \item Herr Solda stellt vor
    \item Das ist keine Sicherheitsunterweisung!
\end{itemize}

\subsection*{Smart Factory B008}
\begin{itemize}
    \item Herr Groß stellt vor
\end{itemize}

\subsection*{Glaskästen}
\begin{itemize}
    \item kurzfristige Änderungen/kranke Profs...
\end{itemize}

\subsection*{PC-Räume}
\begin{itemize}
    \item B101/105 CAD
    \item B108 digitales Werkstofflabor
    \item B107 Informatik
    \item B114 Simulationen
\end{itemize}

\subsection*{Prüfungssekretariat B019}
\begin{itemize}
    \item Fragen zu Prüfungen und Anmeldungen
\end{itemize}

\subsection*{Garten}

\subsection*{Studios B124}
\begin{itemize}
    \item danach den Gang weiter Richtung B129 (Lehrredaktion) und dann runter wegen Klausuren
    \item Stellen sich selbst vor
\end{itemize}

\subsection*{Werkstofflabor B073}
\begin{itemize}
    \item Juliane stellt nur Mittwoch vor
    \item Dienstag nur von außen zeigen
\end{itemize}

\subsection*{Motorsport; E-Labor B064}
\begin{itemize}
    \item Stellen sich selber vor
\end{itemize}

\subsection*{Help C005.1}
\begin{itemize}
    \item Beratungsstelle Familie, Kinder, etc...
    \item Studium und Familie miteinander in Einklang zu bringen
    \item Beratungs- und Vermittlungsangebote
    \item Besonders Schwangerschaft und Geburt
\end{itemize}

\subsection*{Eltern-Kind-Arbeitszimmer C003}
\begin{itemize}
    \item kann man buchen über SIS
\end{itemize}

\subsection*{AStA}
\begin{itemize}
    \item Sportprogramm, eigenes aber auch der Uni Bonn
    \item Man bekommt das Geld über den AStA zurück, das man für den Uni Sport Bonn bezahlen musste
    \item Planen Events: Campus Beben, Sommerfest, Spieleabend
    \item Soziales (/Geld): Darlehen, Rechtsberatung, Windelgeld, Freitisch (kostenlose Mahlzeiten)
\end{itemize}

\subsection*{Bibliothek A105}
\begin{itemize}
    \item Bücher ausleihen
    \item Bücher Kaufempfehlungen
    \item Arbeitsplätze
    \item Buchbare Arbeitsräume
\end{itemize}

\subsection*{Studierwerkstatt A163}
\begin{itemize}
    \item Zeiten hängen aus
    \item Frau Hodel ist immer mal wieder da
    \item Angestellte Studierende helfen auch
    \item Checkt den LEA Kurs ab
\end{itemize}

\subsection*{FB-Sek A157}
\begin{itemize}
    \item organisatorischen Fragen zu Studium, Abschlussarbeiten, Lehrveranstaltungen, Lehraufträgen oder studentischen Hilfskraftstellen
    \item Ansprechen: Prüfungssekretariat B019 bei Prüfungsfragen!
\end{itemize}

\end{document}
 